% ---------------------------------------------------------------------
\section{Problem Formülasyonu}\label{sec:problem}
% ---------------------------------------------------------------------

\subsection{İstasyon, EV'ler ve enerji dinamiği}\label{sec:station}

Bir CPO, tek bir trafonun arkasında $N$ şarj ünitesine $\mathcal N=\{1,\dots,N\}$ sahip bir istasyon işletir; her şarj ünitesine en fazla bir EV bağlanır ve bu EV'ye ait büyüklükler şarj ünitesinin $i$ indisini taşır.
Zaman, $\Delta t$ [h] uzunluğundaki $t=0,\dots,T-1$ adımlarına bölünmüştür.
İstasyon her adımda trafo anma gücünü $P_{\sm{tr}}(t)$, esnek olmayan yükü $P^{\sm{inf}}(t)\ge0$ ve fotovoltaik (PV) üretimi $P^{\sm{pv}}(t)\ge0$ kW cinsinden gözlemler.
Enerji $c^{\sm{ch}}(t)$ fiyatından satın alınır ve $c^{\sm{dis}}(t)$ [\texteuro/kWh] fiyatından satılır; her iki fiyat da negatif olabilir.

EV $i$, $t_i^{\sm{arr}}$ adımında varır ve son kontrol edilen adım olan ayrılış adımı $d_i$'ye kadar kontrol edilir; dolayısıyla ayrılış enerjisi $E_i(d_i+1)$'dir; burada $E_i(t)$, $t$ adımının başındaki batarya enerjisidir.
$t$ adımında EV'nin $\tau_i(t)=d_i-t+1$ kalan eylemi vardır ve $\mathcal N(t)$, $n(t)$ adet dolu şarj ünitesinin kümesidir.
$p_i(t)$ [kW], $i$ şarj ünitesine komut edilen güç olsun; bu güç şarj için pozitif, V2G deşarjı için negatiftir. $\tilde p_i(t)$ ise simülatörün raporladığı güç olsun.
$p^{+}=\max(p,0)$ ve $p^{-}=\max(-p,0)$ olmak üzere batarya enerjisi
\begin{equation}\label{eq:dynamics}
E_i(t+1)=E_i(t)+\Delta t\,g_i\big(p_i(t)\big),\qquad g_i(p)=\eta_i^{\sm c}p^{+}-p^{-}/\eta_i^{\sm d},
\end{equation}
biçiminde değişir; burada $\eta_i^{\sm c},\eta_i^{\sm d}\in(0,1]$ şarj ve deşarj verimleridir.
Adım başına fiziksel limitler gücü
\begin{equation}\label{eq:limits}
\pmn_i(E)\le p_i(t)\le \bar P_i(E),\qquad \pmn_i(E)=-\min\Big\{\Pdis_i,\ \eta_i^{\sm d}\frac{(E-E_i^{\min})^{+}}{\Delta t}\Big\},\qquad E=E_i(t),
\end{equation}
aralığıyla sınırlar; burada en yüksek şarj gücü $\bar P_i(E)\ge0$ şarj gücü azalmasını ve kapasite limitini içerir, $\Pnom_i=\sup_E\bar P_i(E)$'dir ve $\Pdis_i$ en yüksek deşarj gücüdür.
Kuram genel $g_i$ eşlemesi için geçerlidir; ancak EV2Gym, batarya tarafındaki enerji değişimini şarj ünitesi gücü olarak raporlar, faturalandırır ve ölçer \cite{orfanoudakis2025ev2gym}.
Bu nedenle deneylerimizde $\eta_i^{\sm c}=\eta_i^{\sm d}=1$ alınmaktadır; bu kabul, V2G arbitrajındaki gidiş-dönüş kaybını ortadan kaldırmaktadır (\Sec{sec:limitations}).

\subsection{Ayrılış gereksinimi ve trafo limiti}\label{sec:requirement}

Her EV, istenen bir ayrılış enerjisi $E_i^{\sm{des}}$ talep eder ve istasyon düzeyindeki bir talep işleme hattı bu talebi kabul edilen gereksinim $\Ereq_i$'ye dönüştürür.
İşleme hattı, her test tekrarı (replay) için herhangi bir denetleyiciden önce bir kez çalışır; karşılaştırılan her yöntem ve her iki optimum aynı kabul edilen gereksinimi alır.
Varışta talep, kalış süresi boyunca tam güçte şarjın ulaşabileceği enerjiyle, bu değerin yaklaşık $10^{-6}$ kWh altında olacak şekilde sınırlandırılır.
Ardından kabul kontrolü küçük bir olurluluk doğrusal programı (LP) çözer ve bağlı EV'lerin tümüne öngörülen trafo limitleri altında hizmet verilemeyecekse yeni talebi düşürür.
Daha önce kabul edilmiş talepler hiçbir zaman azaltılmaz.
Dolayısıyla
\begin{equation}\label{eq:cap}
\Ereq_i\le\min\big\{E_i^{\sm{des}},\ \text{kalış süresince }E_i^{\sm{arr}}\text{ değerinden tam güçle erişilebilen enerji}\big\}
\end{equation}
olur; ayrıntılar Bölüm~\ref{S-sec:pipeline} içinde verilmektedir.
Ayrılış gereksinimi şudur:
\begin{equation}\label{eq:requirement}
E_i(d_i+1)\ \ge\ \Ereq_i .
\end{equation}
Bu sınırlama olmasaydı bazı varışlara hiçbir denetleyici hizmet veremezdi ve bunlara ilişkin bir garanti anlamsız olurdu.
Her yöntem için eksik enerji, ortak kabul edilen gereksinime göre ölçülür; işleme hattının çıkardığı enerji ise ayrıca raporlanır.

Net trafo gücü $P^{\sm{inf}}(t)-P^{\sm{pv}}(t)+\sum_{i}p_i(t)$, $[-P_{\sm{tr}}(t),P_{\sm{tr}}(t)]$ aralığında kalmalıdır; dolayısıyla EV toplamı için trafo limiti
\begin{equation}\label{eq:transformer}
-\Crev(t)\le\sum_{i\in\mathcal N}p_i(t)\le C(t),\qquad
C(t)=P_{\sm{tr}}(t)-P^{\sm{inf}}(t)+P^{\sm{pv}}(t),\qquad \Crev(t)=P_{\sm{tr}}(t)+P^{\sm{inf}}(t)-P^{\sm{pv}}(t)
\end{equation}
biçimindedir.
İleri yön limiti $C(t)$ şarjı, ters yön limiti $\Crev(t)$ ise V2G ile dışarı verilen gücü sınırlar.
Güvenlik katmanı mevcut adım için ölçülen yük ve PV değerlerini kullanırken politika gelecekteki değerlerin yalnızca gürültülü tahminlerini gözlemler.

\subsection{Çevrim içi kontrol problemi ve optimumlar}\label{sec:online}

CPO her adımda $p(t)\in\R^N$ değerini yalnızca $t$ adımında mevcut olan bilgiyi kullanarak seçer.
Bir $\pi$ denetleyicisinin bir epizottaki \texteuro{} cinsinden kârı
\begin{equation}\label{eq:objective}
J^{\pi}=\sum_{t=0}^{T-1}\Delta t\sum_{i\in\mathcal N(t)}\big(c^{\sm{dis}}(t)\,\tilde p_i^{-}(t)-c^{\sm{ch}}(t)\,\tilde p_i^{+}(t)\big)
\end{equation}
olarak tanımlanır ve denetleyici bunun beklenen değerini \Eqs{eq:dynamics}{eq:limits}, her EV için \Eq{eq:requirement} ve her adımda \Eq{eq:transformer} kısıtları altında en büyüklemelidir.
Her yöntem için $w_E=w_P=10$\,\texteuro/kWh ağırlıklarıyla cezalı kâr
\begin{equation}\label{eq:jpen}
J^{\sm{pen}}(\pi)=J^{\pi}-w_E\sum_i\Delta E_i^{\sm{sh}}-w_P\,\Delta t\sum_t\Big[\Big(\sum_i\tilde p_i(t)-C(t)\Big)^{+}+\Big(-\Crev(t)-\sum_i\tilde p_i(t)\Big)^{+}\Big]
\end{equation}
olarak tanımlanır; burada $\Delta E_i^{\sm{sh}}=(\Ereq_i-E_i(d_i+1))^{+}$ ayrılıştaki eksik enerjidir.

HiGHS \cite{huangfu2018parallelizing} ile çözülen iki tam bilgili LP, epizodun tamamını bilir ve referans olarak kullanılır.
Her ikisi de EV2Gym'in şarj gücü azalmasını dış doğrusal kısıt satırlarıyla modeller; bu satırlar tek adımlık bir geçiş bandı ve \%1'lik bir pay dışında kesindir (Lemma~\ref{S-lem:taperrows}); ayrıca her ikisi de EV2Gym'in faturalandırılmayan 0,01 kWh'lik yuvarlama kazancını hesaba katar.
Yumuşak cezalı optimum $J^{\sm{pen,LP}}$, eksik enerji ve aşırı yük için gevşek değişkenlerle \Eq{eq:jpen} ifadesini en büyükler.
Eksik enerji ortak kabul edilen gereksinime göre ölçüldüğünde, eksik enerji ya da aşırı yük içeren yörüngeler de dâhil olmak üzere simüle edilen her yörünge için $J^{\sm{pen}}(\pi)\le J^{\sm{pen,LP}}$ sağlanır (Önerme~\ref{S-prop:softoracle}).
Dolayısıyla ana ölçüt, şarj ünitesi ve epizot başına \texteuro{} cinsinden mutlak açık olarak karşılaştırılan her yöntem için iyi tanımlıdır:
\begin{equation}\label{eq:gap}
\mathrm{gap}^{\sm{pen}}(\pi)=\frac{J^{\sm{pen,LP}}-J^{\sm{pen}}(\pi)}{N}\ \ge0 .
\end{equation}
Epizot başına göreli bir açık kullanılmamaktadır, çünkü $J^{\sm{pen,LP}}$ çoğu zaman sıfıra yakın ya da negatiftir; ek materyal göreli açıkları yalnızca toplamların oranı olarak raporlamaktadır.
Katı optimum $J^{\sm{LP,ceil}}$, toplam eksik enerjiyi en küçükleyen ilk aşamadan sonra ayrılış gereksinimlerini ve trafo limitini kısıt olarak uygular.
Bu optimum, trafo limitine uyan ve daha fazla eksik enerji bırakmayan her denetleyicinin kârını üstten sınırlar (Önerme~\ref{S-prop:oracle}) ve garantinin maliyeti analizinde kullanılır.

\subsection{Kısmen gözlemlenebilir Markov karar süreci}\label{sec:mdp}

Gelecekteki varışlar, kısa bir pencerenin ötesindeki fiyatlar ve gelecekteki yükler gözlemlenmediğinden çevrim içi problem bir POMDP'dir.
Mevcut gözlem üzerinden ileri beslemeli bir aktör ve kritik ile eylem seçilmektedir.
Politika her şarj ünitesi için bir ham eylem $a_i(t)\in\R$ üretir; bu eylem $P^{\sm{ref}}=11$ kW olmak üzere $x_i(t)=P^{\sm{ref}}a_i(t)$ biçiminde kW'a ölçeklenir ve \Sec{sec:projection} içindeki güvenlik katmanı onu herhangi bir sıkıştırma fonksiyonu kullanmadan komut edilen güce eşler.
Politikanın ödülü adımın gerçekleşen kârıdır; ayrılış gereksinimi ya da trafo limiti için ceza içermez, çünkü güvenlik katmanı her ikisini de uygular.
Eğitimde ödül, getirilerin $N$ ile büyümemesi için $r^{\sm{ref}}=c^{\sm{ref}}P^{\sm{ref}}\Delta t$ ve $c^{\sm{ref}}=0.25$\,\texteuro/kWh olmak üzere $N r^{\sm{ref}}$'ye bölünür; raporlanan tüm sonuçlar ise bölünmemiş kârı kullanır.
Varışlar, fiyatlar, yükler ve PV eylemlerden bağımsız olarak değişir, her EV epizot bitmeden ayrılır ve indirim faktörü $\gamma=0.99$, $\Delta t=15$ dk için yaklaşık bir günlük etkin ufuk sağlar.
